% Source: https://artofproblemsolving.com/community/c3h1441560_cpctc_applications
% Stable identity: aops-topic-1441560-q01
\textbf{(a)} Similar triangles have equal corresponding angles, $\angle A=\angle D$, $\angle B=\angle E$, $\angle C=\angle F$, and proportional corresponding sides:
\[
\dfrac{DE}{AB}=\dfrac{EF}{BC}=\dfrac{FD}{CA}=k .
\]
So corresponding \emph{angles} are always congruent, but sides are only proportional. For example, the right triangles with legs $3,4$ and $6,8$ (hypotenuses $5$ and $10$) are similar with $k=2$, and no pair of corresponding sides is congruent.

\begin{asy}
usepackage("amsmath");
size(6cm);
defaultpen(fontsize(10pt));
pair A=(0,0),B=(4,0),C=(0,3);
pair D=(7,0),Ee=(15,0),F=(7,6);
draw(A--B--C--cycle,linewidth(0.9)); draw(D--Ee--F--cycle,linewidth(0.9));
draw(A+(0.3,0)--A+(0.3,0.3)--A+(0,0.3)); draw(D+(0.4,0)--D+(0.4,0.4)--D+(0,0.4));
label("$A$",A,SW); label("$B$",B,SE); label("$C$",C,N);
label("$D$",D,SW); label("$E$",Ee,SE); label("$F$",F,N);
label("$4$",(A+B)/2,S); label("$3$",(A+C)/2,W); label("$5$",(B+C)/2,NE);
label("$8$",(D+Ee)/2,S); label("$6$",(D+F)/2,W); label("$10$",(Ee+F)/2,NE);
label("\mbox{}",(min(currentpicture,user=true)+max(currentpicture,user=true))/2,black); // keeps the text after the figure black when LaTeX typesets the labels inline
\end{asy}

\textbf{(b)} If $AB=DE$, the ratio is $k=\dfrac{DE}{AB}=1$, so $EF=BC$ and $FD=CA$ as well. By SSS the triangles are congruent, and then all corresponding parts are congruent. (Equivalently, $\angle A=\angle D$, $AB=DE$, $\angle B=\angle E$ gives congruence by ASA.) $\blacksquare$
